\section{The arithmetisation}
\label{sec:air}

Figure~\ref{fig:processor} shows instruction chips connected by buses to a program ROM, memory subsystem and lookup tables. For a chip in the sense of Definition~\ref{def:chip}, the column count is its datapath width and the row count is its activity for an execution. A \emph{bus} is an interconnect on which chips exchange tuples. The lookup argument of Section~\ref{sec:air-logup} acts as the wiring by proving that every tuple driven on a bus is consumed exactly once. The area of the design is the sum over chips of rows $\times$ columns (\emph{trace cells}); Section~\ref{sec:performance} reports it per chip, as a gate-count report would.

\begin{figure}[t]
\centering
\resizebox{\textwidth}{!}{% Ziren as a processor block diagram: units (chips) and buses (interactions).
\begin{tikzpicture}[
  every node/.style={font=\small},
  unit/.style={rectangle, draw=zirenblue!85, rounded corners=2pt, text width=2.0cm, minimum height=0.7cm, align=center, fill=zirenblue!7, inner sep=2pt},
  accel/.style={unit, draw=zirenorange!90, fill=zirenorange!10},
  mem/.style={rectangle, rounded corners=2pt, draw=zirenteal!85, text width=2.7cm, minimum height=0.7cm, align=center, fill=zirenteal!8, inner sep=2pt},
  rom/.style={rectangle, rounded corners=2pt, draw=zirengold!90!black, text width=2.7cm, minimum height=0.7cm, align=center, fill=zirengold!12, inner sep=2pt},
  bus/.style={line width=1.15pt, zirenslate!85},
  programbus/.style={bus, zirenblue!90},
  statebus/.style={bus, zirenteal!90!black},
  memorybus/.style={bus, zirenorange!90!black},
  globalbus/.style={bus, zirenslate!95},
  buslab/.style={font=\small, black!70, right},
  tap/.style={-{Latex[length=1.6mm]}, thin, zirenslate!90},
]
  % instruction units (functional units), pitch 2.3 cm
  \node[unit] (alu)   at (0,0)    {ALU\\\texttt{AddSub}, \texttt{Bitwise}};
  \node[unit] (shift) at (2.3,0)  {Shift\\\texttt{ShiftLeft}\\\texttt{(Imm)}};
  \node[unit] (mul)   at (4.6,0)  {Mul/Div\\\texttt{Mul}, \texttt{DivRem}};
  \node[unit] (ctl)   at (6.9,0)  {Control\\\texttt{Branch}, \texttt{Jump}};
  \node[unit] (ld)    at (9.2,0)  {Load/Store\\\texttt{LoadWord}, \ldots};
  \node[unit] (sys)   at (11.5,0) {Syscall\\\texttt{Syscall}\\\texttt{Instrs}};
  \node[accel, text width=2.5cm] (pre) at (14.9,0) {Accelerators\\Keccak, SHA, EC, \ldots};

  % buses
  \draw[programbus] (-1.2,1.35) -- (16.4,1.35) node[buslab] {program bus};
  \draw[statebus] (-1.2,1.0)  -- (16.4,1.0)  node[buslab] {state bus};
  \draw[memorybus] (-1.2,-1.0) -- (16.4,-1.0) node[buslab] {memory bus};
  \draw[bus] (-1.2,-1.35) -- (16.4,-1.35) node[buslab] {byte / range bus};
  \draw[globalbus] (-1.2,-1.7) -- (16.4,-1.7) node[buslab] {global bus};
  \foreach \u in {alu,shift,mul,ctl,ld,sys,pre} {
    \draw[tap] (\u.north) -- ++(0,0.5);
    \draw[tap] (\u.south) -- ++(0,-0.5);
  }
  \draw[tap] (sys.east) -- node[above=1pt, font=\scriptsize, black!70, align=center] {syscall\\bus} (pre.west);

  % tables and memory units
  \node[rom] (prog) at (1.2,2.5)  {\texttt{Program} ROM\\(preprocessed, in vk)};
  \node[rom, text width=3.4cm] (pv)   at (5.6,2.5)  {Public values\\entry/exit state, digests};
  \node[mem] (bump) at (0.6,-2.9)  {\texttt{MemoryBump}\\register anchors};
  \node[rom] (byte) at (4.0,-2.9)  {\texttt{Byte} / \texttt{Range} tables\\(preprocessed)};
  \node[mem] (loc)  at (7.6,-2.9)  {\texttt{MemoryLocal}\\shard first/last touch};
  \node[mem] (glob) at (11.0,-2.9) {\texttt{Global}\\septic digest};
  \node[mem] (init) at (14.6,-2.9) {\texttt{MemoryGlobal}\\\texttt{Init}/\texttt{Finalize}};

  \draw[tap] (prog.south) -- (1.2,1.35);
  \draw[tap] (pv.south) -- (5.6,1.0);
  \draw[tap] (bump.north) -- (0.6,-1.0);
  \draw[tap] (byte.north) -- (4.0,-1.35);
  \draw[tap] (loc.north) -- (7.6,-1.0);
  \draw[tap] (loc.east) -- (glob.west);
  \draw[tap] (glob.north) -- (11.0,-1.7);
  \draw[tap] (init.north) -- (14.6,-1.7);
\end{tikzpicture}
}
\caption{Each executed instruction occupies one opcode-chip row, and every relation between chips crosses a named bus. The program and lookup tables are preprocessed; public values close the state chain and carry the cross-shard digests.}
\label{fig:processor}
\end{figure}

This section defines the constrained execution relation within the system boundary of Section~\ref{sec:ipcore}. The chip widths determine its per-chip trace area (Appendix~\ref{sec:app-chips}), and the buses connect the chips. One row per instruction, one shared memory argument, one lookup and one zerocheck fix the area measured in Section~\ref{sec:performance}. The guest-accelerator set is configuration rather than fixed design: it determines which precompile chips a shard may contain, and adding one changes the chip set and verifying keys (\S\ref{sec:ipcore-config}).

A shard is proved as a set of chips. A chip is a fixed-width table over $\F$ with polynomial constraints of degree at most three on each row and interactions on named buses. Constraints are local to one row and evaluated at a random point of the multilinear domain (Section~\ref{sec:air-zerocheck}), so relations between rows are expressed as bus interactions. Every relation between chips, or between different rows of one chip, is an interaction. This section describes the chip set, instruction frame, register and memory argument, lookup argument and zerocheck.

\subsection{Chip set}
\label{sec:air-chips}

Appendix~\ref{sec:app-chips} catalogues the chips of the machine measured here, with their widths. Three families dominate the trace area of an Ethereum block: the memory-instruction chips (\texttt{LoadWord}, \texttt{StoreWord}, \texttt{MemoryUnaligned}, \texttt{LoadNarrow}, \texttt{StoreNarrow}), the ALU chips, and the \texttt{Global} chip that carries the cross-shard memory argument. The immediate forms of the ALU instructions (\texttt{ADDIU}, \texttt{ANDI}, \texttt{SLTI}, shifts by a constant, \ldots) have chips separate from the register forms: an immediate row has one register read fewer and a narrower frame, and the split saves roughly a tenth of the width of those rows.


Row counts are rounded up to a multiple of 32 rather than to a power of two; the jagged commitment of Section~\ref{sec:pcs} makes the padding cost proportional to the padding, not to the next power of two. Padding rows have $\mathsf{is\_real} = 0$; every constraint and every interaction multiplicity is gated so that padding rows are vacuous.

\subsection{The instruction frame}
\label{sec:air-frame}

A design organised around a central \emph{CPU chip}, as this system's first generation was, has one row per executed cycle that fetches the instruction, performs the register accesses, chains the program counter and forwards a decoded instruction to an opcode chip over a bus; the opcode chip then spends a second row on the same instruction. The current design removes the central table: each instruction chip embeds an \emph{instruction frame} and performs those duties itself, so one executed instruction is exactly one row of one chip. The name is this paper's, not the architecture's: MIPS32 specifies the delay slot but knows nothing of frames, and this one is unrelated to a stack frame. It is the group of columns and interactions that lets a row stand on its own, and it consists of:

\begin{itemize}
  \item \textbf{Program fetch.} The row sends $(\mathsf{pc}, \mathsf{opcode}, \mathsf{op\_a}, \mathsf{op\_b}, \mathsf{op\_c}, \mathsf{imm\_b}, \mathsf{imm\_c})$ on the \emph{program} bus; the preprocessed \texttt{Program} table receives it with the multiplicity recorded at trace generation. The instruction word therefore comes from the program image committed in the verifying key, and the chip's opcode selector is checked against it.
  \item \textbf{Register accesses.} Reads of $\mathsf{op\_b}$ and $\mathsf{op\_c}$ (unless immediate) and the read-modify-write of $\mathsf{op\_a}$ at the sub-cycle positions $1..4$ of the instruction's tick, through the register variant of the memory argument (Section~\ref{sec:air-memory}). A write to register 0 is committed as zero: the frame carries an $\mathsf{op\_a\_0}$ flag and pins the written word to zero when it is set, so chips bind their result through a $(1-\mathsf{op\_a\_0})$ factor rather than a gate.
  \item \textbf{State bus.} The bus that carries the machine state from one row to the next, which OpenVM calls the execution bus. The row receives the tuple $(\mathsf{shard}, \mathsf{clk}, \mathsf{pc}, \mathsf{next\_pc})$ and sends the tuple $(\mathsf{shard}, \mathsf{clk}+\Delta, \mathsf{next\_pc}, \mathsf{next\_next\_pc})$, where the increment $\Delta$ is five for an ordinary instruction, one per sub-cycle position plus one, and a larger bounded value for a syscall, so that the accelerator's memory accesses fit between the calling instruction and its successor (Section~\ref{sec:isa}); the public-values AIR sends the shard's entry state and receives its exit state. Lemma~\ref{lem:chain} is what this buys: the multiset balances only if the rows form a single chain, which is the shard's execution.
  \item \textbf{Clock.} $\mathsf{clk}$ is witnessed as a 16-bit and a 10-bit limb, both range checked, so that clock differences in the memory argument can be decomposed cheaply.
\end{itemize}

\begin{lemma}[State-bus chain]
\label{lem:chain}
In a valid shard (Definition~\ref{def:shard}) with entry state $s_0$ and exit state $s_N$, the real instruction rows can be ordered $r_1, \dots, r_N$ so that $r_i$ receives $s_{i-1}$ and sends $s_i$, with $\mathsf{clk}(s_i) > \mathsf{clk}(s_{i-1})$; every real instruction row is in the chain, and $N$ is the number of real instruction rows.
\end{lemma}
\begin{proof}
Each real instruction row receives exactly one state tuple and sends exactly one, with a strictly larger clock, and the public-values table sends $s_0$ and receives $s_N$. Balance of the state bus means the multiset of sent tuples equals the multiset of received tuples. Consider the directed graph whose vertices are the distinct state tuples and whose edges are the rows (from received tuple to sent tuple). Every vertex other than $s_0$ and $s_N$ has in-degree equal to out-degree, $s_0$ has out-degree one more than in-degree and $s_N$ the reverse, so the edge multiset decomposes into one $s_0 \to s_N$ path and cycles. A cycle is impossible because the clock strictly increases along every edge. Hence the rows form a single path from $s_0$ to $s_N$.
\end{proof}

Figure~\ref{fig:frame} shows three consecutive rows and the buses that bind them. Three frame variants exist: the I-type frame (an immediate as a 4-byte word, used by memory and branch chips), the R-type frame (two register reads), and a \emph{shamt} frame for shifts by a constant, whose immediate is a single scalar. The frame byte-checks the word written to $\mathsf{op\_a}$; this is the single place where a register value is range checked, and it is what lets memory-instruction rows omit byte checks on the words they move (Section~\ref{sec:air-memory}).

\begin{figure}[!tbp]
\centering
\begin{subfigure}{\textwidth}
\centering
\resizebox{\textwidth}{!}{% Panel (a) of the instruction frame: the state chain and the two shared buses
% across three consecutive cycles.  Deliberately carries NO inner datapath
% equations and NO sub-cycle detail -- those are panel (b) -- so that this
% canvas stays narrow enough for its labels to survive scaling to \textwidth.
\begin{tikzpicture}[
  font=\large,
  unit/.style={rectangle, rounded corners=2pt, draw=zirenblue!85, line width=0.6pt,
               minimum width=3.0cm, minimum height=1.5cm, fill=zirenblue!5, align=center},
  reg/.style={rectangle, rounded corners=1pt, draw=zirenteal!85, line width=0.6pt,
              minimum width=1.5cm, minimum height=1.0cm, fill=zirenteal!8, align=center,
              font=\normalsize},
  pub/.style={rectangle, rounded corners=1pt, draw=zirenslate!80, dashed,
              minimum width=1.5cm, minimum height=1.0cm, fill=zirenlight, align=center,
              font=\normalsize},
  pbus/.style={line width=1.4pt, zirenblue!90},
  mbus/.style={line width=1.4pt, zirenorange!90!black},
  sig/.style={-{Latex[length=1.8mm]}, line width=0.6pt, zirenteal!90!black},
  tap/.style={-{Latex[length=1.6mm]}, line width=0.6pt, black!65},
  lbl/.style={font=\normalsize, black!70},
  hdr/.style={font=\normalsize\bfseries},
]
% ---------- the two shared buses
\draw[pbus] (-2.4,2.0) -- (12.4,2.0)
  node[right, font=\normalsize, text=zirenblue!80!black] {program bus};
\draw[mbus] (-2.4,-1.7) -- (12.4,-1.7)
  node[right, font=\normalsize, text=zirenorange!80!black] {memory bus};

% ---------- three opcode rows
\node[unit] (u1) at (0,0)    {};
\node[unit] (u2) at (5,0)    {};
\node[unit] (u3) at (10,0)   {};
\node[hdr] at ($(u1.north)+(0,-0.42)$) {\texttt{AddSubImm}};
\node[hdr] at ($(u2.north)+(0,-0.42)$) {\texttt{Branch}};
\node[hdr] at ($(u3.north)+(0,-0.42)$) {\texttt{LoadWord}};
\node[lbl] at ($(u1.south)+(0,0.40)$) {\texttt{addiu}};
\node[lbl] at ($(u2.south)+(0,0.40)$) {\texttt{bne}};
\node[lbl] at ($(u3.south)+(0,0.40)$) {\texttt{lw} (delay slot)};

% ---------- each row fetches from the program ROM and touches the memory bus
\foreach \u in {u1,u2,u3} {
  \draw[tap] (\u.north) -- (\u.north |- 0,2.0);
  \draw[tap] (\u.south) -- (\u.south |- 0,-1.7);
}
\node[lbl, anchor=south west] at (-2.4,2.06) {fetch: one receive per row};
\node[lbl, anchor=north west] at (-2.4,-1.76) {register and memory accesses (panel (b))};

% ---------- the state chain
\node[pub] (s0) at (-3.4,0) {entry\\state};
\node[reg] (s1) at (2.5,0)  {state};
\node[reg] (s2) at (7.5,0)  {state};
\node[reg] (s3) at (12.5,0) {state};
\draw[sig] (s0) -- (u1.west);
\draw[sig] (u1.east) -- (s1);
\draw[sig] (s1) -- (u2.west);
\draw[sig] (u2.east) -- (s2);
\draw[sig] (s2) -- (u3.west);
\draw[sig] (u3.east) -- (s3);
\node[lbl, anchor=north, align=center] at ($(s1.south)+(0,-0.12)$) {$\mathsf{clk}{+}5$};
\node[lbl, anchor=north, align=center] at ($(s2.south)+(0,-0.12)$) {$\mathsf{clk}{+}10$};
\node[lbl, anchor=north, align=center] at ($(s3.south)+(0,-0.12)$) {$\mathsf{clk}{+}15$};
\node[lbl, anchor=north, align=center] at ($(s0.south)+(0,-0.12)$) {$\mathsf{clk}$};
\node[lbl, anchor=west, align=left] at ($(s3.east)+(0.15,0)$) {to the row\\at $\mathsf{target}$};
\end{tikzpicture}
}
\caption{The state chain across three consecutive cycles. Each row receives the machine state of its cycle from the register on its left and sends that of the next to the register on its right, fetches its instruction from the preprocessed \texttt{Program} ROM on the program bus, and places its accesses on the memory bus.}
\label{fig:frame-flow}
\end{subfigure}

\vspace{0.7em}
\begin{subfigure}{0.76\textwidth}
\centering
\resizebox{\textwidth}{!}{% Panel (b) of the instruction frame: ONE row enlarged, showing the four fixed
% sub-cycle access positions.  A narrow canvas, so it scales up rather than down
% and its pin labels are the largest text in the figure.  The pins sit ON the
% row's lower edge with their labels BELOW it, so they cannot collide with the
% datapath lines inside the box.
\begin{tikzpicture}[
  font=\large,
  chip/.style={rectangle, rounded corners=2pt, draw=zirenblue!85, line width=0.7pt,
               minimum width=7.2cm, minimum height=1.9cm, fill=zirenblue!5},
  mbus/.style={line width=1.6pt, zirenorange!90!black},
  pbus/.style={line width=1.6pt, zirenblue!90},
  wire/.style={line width=0.5pt, black!80},
  sig/.style={-{Latex[length=1.8mm]}, line width=0.7pt, zirenteal!90!black},
  pin/.style={circle, fill=black!80, inner sep=0pt, minimum size=3.4pt},
  npin/.style={circle, draw=black!45, fill=white, inner sep=0pt, minimum size=3.4pt},
  lbl/.style={font=\normalsize, black!70},
  hdr/.style={font=\large\bfseries},
]
% ---------- the row itself (spans y = -0.95 .. 0.95)
\node[chip] (u) at (0,0) {};
\node[hdr, anchor=north west] at ($(u.north west)+(0.14,-0.10)$) {\texttt{LoadWord} row};
\node[lbl, anchor=north east] at ($(u.north east)+(-0.14,-0.12)$) {\texttt{lw \$t3,0(\$t4)}};
\node[lbl, align=left, anchor=south west] at ($(u.south west)+(0.18,0.12)$)
  {$\mathsf{addr} \leftarrow \mathsf{op\_b} + \mathsf{imm}$,\quad
   $\mathsf{op\_a} \leftarrow \mathsf{mem}[\mathsf{addr}]$};

% ---------- fetch, upward to the program bus
\draw[pbus] (-4.6,2.0) -- (4.6,2.0)
  node[right, font=\normalsize, text=zirenblue!80!black] {program bus};
\node[pin] (uf) at (u.north) {};
\draw[sig] (uf) -- (u.north |- 0,2.0);
\node[lbl, anchor=south east] at ($(uf)+(-0.1,0.42)$) {fetch $(\mathsf{pc},\mathsf{op})$};

% ---------- the four access positions, pinned to the lower edge
\draw[mbus] (-4.6,-3.1) -- (4.6,-3.1)
  node[right, font=\normalsize, text=zirenorange!80!black] {memory bus};
\foreach \k/\t in {1/{rd $\mathsf{op\_b}$},2/{},3/{mem rd},4/{wr $\mathsf{op\_a}$}} {
  \pgfmathsetmacro{\px}{-2.7+1.8*(\k-1)}
  \ifx\t\empty
    \node[npin] (p\k) at (\px,-0.95) {};
    \node[lbl, anchor=north, align=center, black!40] at (\px,-1.10) {@\k\\ unused};
  \else
    \node[pin] (p\k) at (\px,-0.95) {};
    \node[lbl, anchor=north, align=center] at (\px,-1.10) {@\k\\ \t};
    \draw[sig] (\px,-2.05) -- (\px,-3.1);
  \fi
}

% ---------- the sub-cycle ruler under the bus
\draw[wire] (-3.3,-3.65) -- (3.3,-3.65);
\foreach \k in {1,...,4} {
  \pgfmathsetmacro{\px}{-2.7+1.8*(\k-1)}
  \draw[wire] (\px,-3.55) -- (\px,-3.75);
  \node[lbl, anchor=north] at (\px,-3.79) {$\mathsf{clk}{+}\k$};
}
\node[lbl, anchor=east, align=right] at (-3.4,-3.65) {access\\positions};
\end{tikzpicture}
}
\caption{One row enlarged: the four fixed sub-cycle access positions. Filled pins are accesses the opcode uses, hollow pins are unused.}
\label{fig:frame-chip}
\end{subfigure}
\caption{The instruction frame binds consecutive opcode rows into one state chain (panel~\subref{fig:frame-flow}) while assigning register and memory accesses to fixed sub-cycle positions (panel~\subref{fig:frame-chip}). The state tuple is $(\mathsf{shard}, \mathsf{clk}, \mathsf{pc}, \mathsf{next\_pc})$; the program tuple received from the ROM is $(\mathsf{pc}, \mathsf{opcode}, \mathsf{op\_a}, \mathsf{op\_b}, \mathsf{op\_c}, \mathsf{imm}, \ldots)$; each memory access sends the previous tuple $(\mathsf{shard},\, \mathsf{clk}{+}\mathsf{pos},\, \mathsf{addr},\, \mathsf{value})$ and receives the new one. In panel~\subref{fig:frame-flow} the \texttt{LoadWord} row sits in the branch's delay slot, so its successor arrives through the state register as the branch target rather than $\mathsf{pc}{+}4$. The rows compute $\mathsf{op\_a} \leftarrow \mathsf{op\_b} + \mathsf{imm}$ with a byte-checked result and witnessed carries (\texttt{AddSubImm}), $\mathsf{taken} \leftarrow [\mathsf{op\_b} \neq \mathsf{op\_c}]$ with $\mathsf{next\_next\_pc} \leftarrow \mathsf{taken}\,?\,\mathsf{target} : \mathsf{next\_pc}{+}4$ (\texttt{Branch}), and the aligned four-lane load of panel~\subref{fig:frame-chip} (\texttt{LoadWord}).}
\label{fig:frame}
\end{figure}

\subsection{Registers and memory}
\label{sec:air-memory}

Registers and memory share one offline memory-checking argument~\citep{blum1994} in the multiset style used by SP1~\citep{sp1}: the state of address $a$ is the tuple $(\mathsf{shard}, \mathsf{clk}, a, \mathsf{value})$ of its last access. An access at time $(s, c)$ \emph{sends} the previous tuple $(s', c', a, v')$ and \emph{receives} the new tuple $(s, c, a, v)$ on the \emph{memory} bus, with $v = v'$ for a read. The multiset balances if and only if every access consumes exactly the tuple produced by the previous access to the same address, provided timestamps strictly increase along each address's chain; the AIR of every access enforces $(s', c') < (s, c)$ by witnessing $c - c' - 1$, or $s - s' - 1$ when the shards differ, in range-checked limbs that cover a $\mathsf{clk} + \mathsf{pos}$ timestamp under a 25-bit clock fence, as a 16-bit limb and a 10-bit high limb (a one-bit \texttt{compare\_clk} selector chooses the branch). Words are four byte columns; the argument carries them as such.
Figure~\ref{fig:memchain} follows one address through two shards: which of its
tuples the memory bus balances, which are left to the cross-shard digest, and
where the chain begins and ends.

\begin{lemma}[Memory argument]
\label{lem:memory}
In a valid shard, for every address $a$ the accesses to $a$ form a chain ordered by $(\mathsf{shard}, \mathsf{clk})$ in which each access consumes exactly the tuple produced by the previous access to $a$, starting from the tuple supplied by the initial-state row of $a$ and ending at the tuple consumed by its final-state row; in particular every read returns the value of the last write.
\end{lemma}
\begin{proof}
Each access sends its previous tuple and receives its new tuple on the memory bus, and its constraints enforce that the new timestamp strictly exceeds the previous one. Balance of the memory bus, restricted to the tuples with address $a$, gives a graph as in Lemma~\ref{lem:chain} with the initial-state row as source and the final-state row as sink; strictly increasing timestamps exclude cycles, so the accesses form one path, and along the path a read carries its previous value forward unchanged.
\end{proof}

\paragraph{Registers.} The \texttt{MemoryBump} chip inserts, for every register at its first touch in a shard, a shadow read at $(\mathsf{shard}, 0)$. Since $\mathsf{clk}$ restarts at zero in every shard and real accesses sit at positions $1..4$ of a tick, the previous access of a register is always in the current shard. The register variant of the access columns therefore drops $\mathsf{prev\_shard}$, the compare selector and one limb, which is re-derived linearly from the others: 10 columns become 7 on every register access of every instruction.

\paragraph{Shard-local memory.} Within a shard, the first access to a word must consume a tuple that no access in the shard produced, and the last access produces a tuple that nothing in the shard consumes. The \texttt{MemoryLocal} chip records, for every word touched in the shard, its initial and final $(\mathsf{shard}, \mathsf{clk}, \mathsf{value})$; it sends the initial tuple on the memory bus (so the first access can consume it) and receives the final one, and it forwards both as \emph{global} interactions (Section~\ref{sec:air-global}) so that they are matched across shards. On a \texttt{reth} block (the Ethereum execution-client workload of Section~\ref{sec:performance}) a few hundred thousand words are touched per shard, most of them the hot working set touched again in the next shard; the two global rows per touched word are the largest single contribution to trace area, which is why the shard size matters (Section~\ref{sec:performance}).

\paragraph{Whole-execution memory.} \texttt{MemoryGlobalInit} carries the memory's initial contents (the program image and the hint data written into untouched words) and \texttt{MemoryGlobalFinalize} its final contents. Both witness the word as 32 boolean columns, which is what makes every value that enters memory byte-shaped by construction; addresses are 32-bit decomposed and proved strictly increasing along the table through a running comparison against the previous address, with the address cursors published as public values so the recursion can check that consecutive tables tile the address space without overlap. The program image itself is bound by the verifying key: its digest (Section~\ref{sec:air-global}) is a public parameter, and the init table must consume it.

\paragraph{What is byte-checked.} Every value that enters a register is byte-checked by the frame at the write, and every value a precompile writes to memory is byte-checked or bit-decomposed by the precompile chip; memory-instruction rows therefore omit the two additional byte lookups per word that accesses otherwise carry.

\begin{figure}[!tbp]
\centering
\resizebox{\textwidth}{!}{% The memory argument for ONE address: a chain that is shard-local on the
% memory bus and closed across shards only by the septic digest.  Solid edges
% are memory-bus tuples, dashed edges are global interactions.  The two shards
% are stacked rather than placed side by side, so that the canvas stays close
% to 2:1 and the tuple labels survive scaling to \textwidth.
\begin{tikzpicture}[
  font=\large,
  acc/.style={rectangle, rounded corners=2pt, draw=zirenblue!85, fill=zirenblue!6,
              line width=0.7pt, minimum width=2.2cm, minimum height=1.25cm,
              align=center, font=\normalsize},
  loc/.style={rectangle, rounded corners=2pt, draw=zirenteal!85, fill=zirenteal!9,
              line width=0.7pt, minimum width=2.2cm, minimum height=1.25cm,
              align=center, font=\normalsize},
  glob/.style={rectangle, rounded corners=2pt, draw=zirenslate!80, fill=zirenlight,
               line width=0.7pt, minimum width=2.4cm, minimum height=1.25cm,
               align=center, font=\normalsize},
  bus/.style={-{Latex[length=2.4mm]}, line width=1.2pt, zirenorange!85!black},
  gl/.style={-{Latex[length=2.4mm]}, line width=1.1pt, zirenslate!85, dashed,
             dash pattern=on 3.5pt off 2.5pt},
  tup/.style={font=\normalsize, text=zirenorange!70!black, inner sep=1.5pt},
  gtup/.style={font=\normalsize, text=zirenslate, align=center, inner sep=2.5pt,
               fill=white},
  note/.style={font=\small, black!70, align=center, inner sep=1.5pt},
  hdr/.style={font=\normalsize\bfseries, black!65},
  shardbox/.style={rectangle, rounded corners=3pt, draw=black!35, dashed,
                   dash pattern=on 3pt off 2.5pt, line width=0.8pt},
]

% ================= shard s: the memory bus balances inside this box
\draw[shardbox] (-1.35,-1.8) rectangle (13.35,1.75);
\node[hdr, anchor=north west] at (-1.2,1.66) {shard $s$};

\node[loc] (li) at (0,0)    {\texttt{MemoryLocal}\\initial};
\node[acc] (r1) at (4.4,0)  {\texttt{lw}\\read at $c_1$};
\node[acc] (w1) at (8.8,0)  {\texttt{sw}\\write at $c_2$};
\node[loc] (lf) at (13.2,0) {\texttt{MemoryLocal}\\final};

\draw[bus] (li) -- (r1);
\draw[bus] (r1) -- (w1);
\draw[bus] (w1) -- (lf);
\node[tup] at (2.2,1.02)  {$(s,\,0,\,a,\,v_0)$};
\node[tup] at (6.6,1.02)  {$(s,\,c_1,\,a,\,v_0)$};
\node[tup] at (11.0,1.02) {$(s,\,c_2,\,a,\,v_1)$};
\node[note] at (6.6,-1.08) {a read carries $v_0$ forward};
\node[note] at (11.0,-1.08) {a write installs $v_1$};

% ================= shard s+1: the chain resumes at the tuple shard s left
\draw[shardbox] (-1.35,-6.0) rectangle (9.05,-2.45);
\node[hdr, anchor=north east] at (8.9,-2.54) {shard $s{+}1$};

\node[loc] (li2) at (0,-4.2)   {\texttt{MemoryLocal}\\initial};
\node[acc] (r2)  at (4.4,-4.2) {\texttt{lw}\\read at $c_1'$};
\node[loc] (lf2) at (8.8,-4.2) {\texttt{MemoryLocal}\\final};

\draw[bus] (li2) -- (r2);
\draw[bus] (r2) -- (lf2);
\node[tup] at (2.2,-3.18) {$(s,\,c_2,\,a,\,v_1)$};
\node[tup] at (6.6,-3.18) {$(s{+}1,\,c_1',\,a,\,v_1)$};

% ================= the same tuple, forwarded as two global interactions
\draw[gl] (lf.south) -- (13.2,-2.13) -- (0,-2.13) -- (li2.north);
\node[gtup] at (7.0,-2.13)
  {one tuple: sent in shard $s$, received in shard $s{+}1$, cancels in $D$};

% ================= the whole-execution tables at the two ends of the chain
\node[glob] (gi) at (-4.6,0)    {\texttt{MemoryGlobalInit}};
\node[glob] (gf) at (12.7,-4.2) {\texttt{MemoryGlobalFinalize}};
\draw[gl] (gi) -- (li);
\draw[gl] (lf2) -- (gf);
\node[note, text width=4.0cm] at (-4.6,-1.35)
  {program image and hint data; its digest is pinned by $\mathsf{vk}$};
\node[note, text width=4.2cm] at (12.7,-5.55)
  {addresses strictly increasing; the cursors are public values};

% ================= legend
\node[note, anchor=north west, align=left] at (-5.9,-6.55)
  {\textcolor{zirenorange!85!black}{\rule[0.22em]{1.4em}{1.1pt}}\ \
   memory bus: balances within one shard, and $(s,c)$ strictly increases
   along the chain \qquad
   \textcolor{zirenslate!85}{\rule[0.22em]{0.45em}{1.1pt}\hspace{0.18em}%
   \rule[0.22em]{0.45em}{1.1pt}}\ \
   global interaction: summed into the digest $D$, cancels only over the whole
   execution};
\end{tikzpicture}
}
\caption{The chain of one address $a$ across a shard boundary. A shard proof can balance only the tuples whose two sides it contains, so the \texttt{MemoryLocal} rows stand in for the accesses it cannot see: the initial row supplies the tuple the shard's first access consumes, and the final row consumes the tuple its last access produces. Those two tuples are then the \emph{same} tuple on either side of the boundary, carried with opposite signs into the digest $D$ of Section~\ref{sec:air-global}, which is why the chain closes over the execution without any shard proof depending on another. The two ends belong to the whole-execution tables rather than to a shard.}
\label{fig:memchain}
\end{figure}

\subsection{Cross-shard interactions: the global digest}
\label{sec:air-global}

Interactions whose two sides live in different shards cannot be balanced by a per-shard fraction sum. These include the initial and final tuples of every shard-touched word, memory initialisation and finalisation, and syscall pairings. \Zirenver{} carries them in an elliptic-curve digest of the kind SP1 introduced~\citep{sp1}. An interaction tuple $\bm m \in \F^7$ is mapped to a point of $E: y^2 = x^3 + 3\zeta x - 3$ over the degree-seven extension $\F_{p^7} = \F[\zeta]/(\zeta^7 + 2\zeta - 8)$; sends contribute the point and receives its negation, and the sum over the whole execution must be the identity. A shard proof publishes its running sum, and the block verifier adds them.

The \texttt{Global} chip has one row per such interaction. It witnesses the point, checks the curve equation, and fixes the sign of $y$ by confining its top limb to one of two disjoint bands, one for sends and one for receives. The running sum is carried on an accumulation bus, row $i$ receiving $(i, S_i)$ and sending $(i+1, S_{i+1})$ with the chord identities enforcing the addition, and the public-values table closes the chain.

The two exceptional cases of the chord law are excluded separately, and by different mechanisms. Both identities carry the factor $(x_2 - x_1)$:
\[
  C_x = (x_1 + x_2 + x_3)(x_2 - x_1)^2 - (y_2 - y_1)^2, \qquad
  C_y = (y_1 + y_3)(x_2 - x_1) - (y_2 - y_1)(x_1 - x_3).
\]
At $P_2 = -P_1$ we have $C_x = -4y_1^2$, which is non-zero whenever $y_1 \neq 0$, so the identity case is rejected by $C_x$ itself. At $P_2 = P_1$ both coordinate differences vanish, so $C_x$ and $C_y$ hold for \emph{every} $P_3$ and the identities alone would leave the next running sum unconstrained but for the curve equation. The chip therefore witnesses $(x_2 - x_1)^{-1}$ and constrains $(x_2 - x_1)\cdot(x_2-x_1)^{-1} = 1$ on every real row, which makes $x_2 \neq x_1$ a constraint rather than an assumption; since both points are on the curve, $x_2 = x_1$ forces $y_2 \in \{y_1, -y_1\}$ and there is no third case. The first row is separately safe: the sum starts at a fixed public offset point $D$ chosen with $\pm D$ outside the image of the map. What remains is a running sum at some later row coinciding with the next interaction's point, which the inverse makes unprovable rather than unconstrained, and whose probability Assumption~\ref{ass:digest} bounds by the unquantified $\varepsilon_{\mathrm{digest}}$.

One departure from SP1 matters. SP1 hashes the message to the curve with a permutation constrained inside the chip, so its points are random-oracle outputs. \Zirenver{} applies no hash: the message \emph{is} the $x$-coordinate, up to a prover-chosen byte $\mathsf{offset}$ in the low positions of $x_6$ that makes the curve's right-hand side a square, so that a point with that $x$ exists. That is what reduces the chip from a few hundred columns to a few dozen, and it changes the security argument rather than merely its constants.

\subsection{Security of the digest}
\label{sec:air-digest-security}

The lift is a map on \emph{pairs}, not on messages: $\mathsf{enc}(\bm m, \mathsf{offset})$. The offset is a witness column, and the chip constrains only that it is a byte occupying the low positions of $x_6$; it does not constrain it to be canonical. The honest prover takes the smallest admissible offset. The implementation scans $\mathsf{offset} = 0, \dots, 255$ and stops at the first value whose $x$ makes $x^3 + 3\zeta x - 3$ a square in $\F_{p^7}$, so that the point exists. Minimality is not an AIR constraint, so a satisfying trace may carry any admissible offset. We therefore write $\mathcal P$ for the \emph{reachable pairs} $(\bm m, \mathsf{offset})$ that an interaction of a satisfying trace can carry, given the range facts each sending chip enforces and the byte bound on the offset, and reserve $\mathcal M$ for the projection of $\mathcal P$ onto messages. $\mathcal M$ is far smaller than $\F^7$, since every limb of a memory or syscall tuple is a byte, a bounded clock limb or a bounded address, and the whole argument rests on that; the offset multiplies the reachable set by at most $256$.

Two consequences follow from the offset being free. First, cancellation is over pairs: a send and its matching receive contribute $\pm\mathsf{enc}(\bm m, \mathsf{offset})$ and cancel only if they agree on the offset. A trace whose two sides choose different admissible offsets for the same message does not balance and is rejected; the freedom costs completeness, not soundness, on that axis. Second, and this is the axis that matters below, the adversary searching for a vanishing combination searches over $\mathcal P$ rather than $\mathcal M$.

\begin{lemma}[injectivity of the lift]
\label{lem:digest}
Assume the message limbs satisfy the range facts the sending chips enforce, so that $256\,m_6 + \mathsf{offset} < p$. Then $\mathsf{enc}$ is injective: limbs $x_0, \dots, x_5$ are $m_0, \dots, m_5$; $x_6$ is $m_6$ shifted above a low byte holding the offset, which the bound makes separable; and the two sign bands are disjoint. So $\bm m$ and the offset are recoverable from the point.
\end{lemma}

Injectivity makes the digest a faithful encoding of the multiset of \emph{pairs}, but it is not what makes the argument sound. If a claimed execution has send and receive multisets $S \neq R$ over $\mathcal P$ whose digests agree, then $\sum_i c_i\,\mathsf{enc}(\bm m_i, \mathsf{offset}_i) = O$ for the non-zero integer vector $c = \mathrm{mult}_S - \mathrm{mult}_R$. Equality of encoded multisets implies equality of message multisets only after projecting away the offsets, which is legitimate precisely because the lift is injective on pairs (Lemma~\ref{lem:digest}). Soundness is therefore exactly:

\begin{assumption}[no reachable vanishing combination]
\label{ass:digest}
Let $\lambda$ be the security parameter, let $N_{\mathrm{int}}$ bound the interactions a block proof may contain and $B$ the per-interaction multiplicity bound, and let $\mathcal P$ be the reachable pairs above. For every adversary $\mathcal A$ running in time at most $T(\lambda)$ and making at most $q(\lambda)$ queries to the hash functions of the protocol,
\[
  \Pr\big[\, \mathcal E(\mathcal A) \,\big] \;\le\; \varepsilon_{\mathrm{digest}}(\lambda, T, q, N_{\mathrm{int}}, B),
\]
where $\mathcal E(\mathcal A)$ is the event that $\mathcal A$ outputs reachable pairs $\{(\bm m_i, \mathsf{offset}_i)\}_{i \le N_{\mathrm{int}}} \subseteq \mathcal P$ together with a non-zero integer vector $c \in \mathbb{Z}^{N_{\mathrm{int}}}$ satisfying $\|c\|_\infty \le B$ and $\sum_i c_i\, \mathsf{enc}(\bm m_i, \mathsf{offset}_i) = O$ in $E(\F_{p^7})$.
\end{assumption}
The assumption is that $\varepsilon_{\mathrm{digest}}$ is at most the level Section~\ref{sec:iopp-soundness} targets, at the interaction and multiplicity bounds of a block proof; that is what Theorem~\ref{thm:composition} takes as a hypothesis, and this machine's parameters are fixed rather than a family, so the statement is concrete. We do not claim a value for it; the paragraph below explains why, and Section~\ref{sec:iopp-soundness} carries it as a named open component rather than folding it into a bit total.

This is the weakest link in the accounting, and we state it rather than a bit figure for two reasons. Hashing to the curve, as SP1 does, would make the points random-oracle outputs and support a clean discrete-logarithm statement about points no adversary can steer. Here the prover chooses points from a structured set, and reads a message off any target $x$-coordinate by inverting the limb layout of Lemma~\ref{lem:digest}; the uncanonicalised offset widens that set by a further factor of at most $256$. The sparsity of $\mathcal P$ blocks the obvious attack of choosing a generator $G$, selecting scalars that sum to zero and reading their $x$-coordinates as messages, since almost no such $x$ has the limb shape a reachable pair requires. This reasoning depends on every sending chip bounding every contributed limb, an arithmetisation obligation not collected in one certificate today. Sparsity also does not yield a bit estimate: the relevant generalised-birthday cost falls with the number of lists and must be bounded using the interactions a prover can place in a trace, not only $|E(\F_{p^7})| \approx 2^{217}$. That analysis remains open. Section~\ref{sec:iopp-soundness} therefore lists the digest as a separate computational component.

\subsection{Trace uniqueness}
\label{sec:air-unique}

\begin{theorem}[From chip determinism to trace uniqueness]
\label{thm:unique}
Fix a program image, public values and a valid shard. If every chip of the shard is deterministic in the sense of Definition~\ref{def:det}, then the rows reachable from the entry state along the state and memory buses are determined by the program image and the public values: any two valid shards with the same program image and public values agree on those rows up to permutation.
\end{theorem}
\begin{proof}
By Lemma~\ref{lem:chain} the instruction rows of either shard form one chain from the entry state; write $r_1, \dots, r_N$ for the first shard's chain. We show by induction on $i$ that $r_i$ has the same inputs and outputs in both shards. The inputs of $r_i$ are its received state tuple $s_{i-1}$, its fetched instruction, and the previous values of the addresses it accesses. $s_0$ is public and $s_{i-1}$ is the output of $r_{i-1}$ by induction. The fetched instruction is a function of $\mathsf{pc}(s_{i-1})$ because the program bus is received only by the preprocessed program table, which is committed. The previous value of an address is, by Lemma~\ref{lem:memory}, the value written by the last earlier access to that address, or the initial value from the memory-initialisation table, which is bound to the program image and the public address cursors; earlier accesses belong to rows $r_j$, $j < i$, whose outputs agree by induction. So the inputs of $r_i$ agree, and determinism of its chip gives that its outputs agree. Rows of non-instruction chips (the memory-initialisation and finalisation tables, the shadow reads, the digest rows, the tables) are determined in the same way from the bus tuples they receive, which are outputs of instruction rows or public. Rows that participate in no bus reachable from the entry state are not covered by this argument; the arithmetisation admits none, because every real row carries a frame or a memory access, but that is a property of the chip set rather than of the buses, and we do not prove it here.
\end{proof}

The theorem is stated, not machine-checked; its premises are the chip-level determinism statements, of which \ndetproved{} of \ndettotal{} are machine-checked today and the rest are assumed (Section~\ref{sec:verification-det}), and the two bus lemmas are the standard arguments of the multiset construction.

\subsection{Preprocessed tables}

The \texttt{Byte} table has $2^{16}$ rows indexed by a byte pair and one multiplicity column per operation (bitwise operations, shifts, comparisons and 8/16-bit range checks); every carry, borrow, comparison and range check in the machine is an interaction with it. The \texttt{Range} table serves the clock limbs, and the \texttt{Program} table holds the decoded program image. Preprocessed tables are committed at setup as part of the verifying key and opened at the same points as the main trace.

\subsection{The shard-level lookup argument}
\label{sec:air-logup}

This is the bus argument of \S\ref{sec:argument-components}.

One \LogUp{} instance proves all interactions of a shard on the program, state, memory, byte, range, syscall, global and accumulation buses. After committing the trace, the verifier samples $\alpha, \beta \in \EF$. Every interaction $(\mathsf{bus}, \bm f, m)$ contributes $m / (\alpha - \mathsf{bus} - \sum_{j \ge 1} \beta^{j} f_j)$, and the prover shows that sends and receives have equal sums. These fractions are leaves of a layered GKR circuit. Sumcheck proves its transitions from the root down, and the final layer produces multilinear evaluation claims on the participating columns for the polynomial commitment. A reth shard has on the order of $10^8$ leaves. Section~\ref{sec:performance} shows that this circuit is the largest GPU-kernel consumer; its leaf count is the number of (row, interaction) pairs, so interactions per row are a separate cost from columns.

Soundness of the fraction identity is $\deg/|\EF|$ per sumcheck round plus the collision probability of the fingerprint, which grows with the number of leaves; Section~\ref{sec:iopp-soundness} lists it as one component of the accounting. The checked-in schedule grinds the \LogUp{} challenge by 16 bits (Table~\ref{tab:whir-schedule}); the performance baseline of Section~\ref{sec:perf-changes} was measured with none (\S\ref{sec:iopp-soundness}).

\subsection{The shard-level zerocheck}
\label{sec:air-zerocheck}

This is the constraint argument of \S\ref{sec:argument-components}.

The polynomial constraints of all chips are proved by \emph{one} zerocheck. The verifier samples $\tau$ for a random linear combination across chips and $\mu$ across the constraints of a chip; for each chip the polynomial $C_c(\bm X) = \sum_k \mu^k\, c_{c,k}(\text{row}(\bm X))$ must vanish on the chip's real rows. The prover proves $\sum_{\bm b} \eqf(\bm b, \bm r)\, \sum_c \tau^c C_c(\bm b) = 0$ for a random $\bm r$ by sumcheck over the tallest chip's variables, with shorter chips handled by an analytic ``virtual $\ge$'' factor rather than padding, and the final evaluation claims on the chips' columns are again handed to the commitment. The constraint degree is three, which fixes the degree of the round polynomials at four; the constraint-evaluation kernel on the GPU is compiled per chip from the same constraint description that the verifier uses, except for the widest accelerator chips, which are interpreted (Table~\ref{tab:census}), and a mismatch between the two is caught by host verification (Section~\ref{sec:performance}).

The zerocheck's claims are seeded from the \LogUp{} openings so that the two protocols share one opening of each column: after the two sumchecks, what remains is a set of (chip, column, point, value) claims, of the order of $3.6 \cdot 10^4$ per shard, and turning those into a single claim on a single committed polynomial is the job of the commitment layer.
